%! Advanced Factoring — Unit 3, Lesson 3.2 — Warm-Up (Answer Key)
\documentclass[10pt]{article}
\usepackage{algebra2-article}
\usepackage{algebra2-key}   % pulls in -boxes; do NOT also load -boxes
\newcommand{\TallMath}[1]{$\displaystyle #1\rule[-1.4em]{0pt}{3.2em}$}

\begin{document}

\pageheader{Unit 3, Lesson 3.2}{Warm-Up --- Answer Key}
\namedateperiod

\vspace{0.10in}

\begin{notesbox}{Getting Ready: Skills We'll Use Today}
\small Yesterday you multiplied polynomials \emph{forward}. Today you will run them \emph{backward}.
These four skills are what that takes.

\vspace{5pt}
\textbf{1. Greatest common factor.} The GCF of two monomials uses the largest common number and the
\emph{lowest} power of each shared variable.

\vspace{2pt}
GCF of $18$ and $24$: \ans{6} \qquad GCF of $12x^3$ and $18x^5$: \ans{$6x^3$} \qquad
GCF of $10x^2y$ and $15xy^3$: \ans{$5xy$}

\vspace{3pt}
Factor it out: \; $10x^2 + 15x = $ \ans{$5x(2x+3)$}

\vspace{6pt}
\textbf{2. Factoring you already own} (from Lesson 2.2).

\vspace{2pt}
$x^2 + 7x + 12 = $ \ans{$(x+3)(x+4)$} \qquad $x^2 - 9 = $ \ans{$(x-3)(x+3)$}

\vspace{3pt}
Which one used the \emph{difference of squares} pattern $a^2-b^2=(a+b)(a-b)$? \ans{the second}

\vspace{6pt}
\textbf{3. A product from yesterday --- watch what survives.} Multiply, then collect like terms.

\vspace{2pt}
$(x + 2)(x^2 - 2x + 4)$

\vspace{2pt}
Six products: \ans{$x^3$} \, \ans{$-2x^2$} \, \ans{$4x$} \, \ans{$2x^2$} \, \ans{$-4x$} \,
\ans{$8$}

\vspace{3pt}
Collected answer: \ans{$x^3+8$} \quad How many of the six terms survived? \ans{2}

\vspace{6pt}
\textbf{4. Perfect cubes.} Fill the table --- you will need these on sight today.

\vspace{2pt}
\begin{center}
\renewcommand{\arraystretch}{1.4}
\begin{tabular}{|c|c|c|c|c|c|}
\hline
$n$ & $1$ & $2$ & $3$ & $4$ & $5$ \\ \hline
$n^3$ & \ans{1} & \ans{8} & \ans{27} & \ans{64} & \ans{125} \\ \hline
\end{tabular}
\end{center}

\vspace{2pt}
Also: $(2x)^3 = $ \ans{$8x^3$} \qquad $(3y)^3 = $ \ans{$27y^3$}

\end{notesbox}

\begin{teachernote}
Item~3 is the whole point of this Warm-Up: students \emph{derive} today's new identity before it is
named. Debrief it aloud --- ask why four of the six terms vanished ($-2x^2$ with $+2x^2$, and $+4x$
with $-4x$) --- because that pairwise cancellation is exactly why the sum of cubes has such a short
answer. When you reach the cube identities in the notes, point back at this line: $(x+2)(x^2-2x+4)
= x^3+8$ \emph{is} $a^3+b^3=(a+b)(a^2-ab+b^2)$ with $a=x$, $b=2$. Item~1's last two answers guard the
year's most common GCF error --- taking the \emph{highest} shared power instead of the lowest.
Item~4 takes ninety seconds but the cube identities are unusable without it.
\end{teachernote}

\end{document}
